\vspace{2.5pt plus 3pt minus 0.5pt}\noindent\hspace*{1.2em}\kbvnum{(10)}\hspace{0.55em}ಈ ಅಂಶಗಳನ್ನೊಳಗೊಂಡಂತೆ ಅರ್ಥವತ್ತಾದ ವಾಕ್ಯಗಳನ್ನು ರಚಿಸಿ :\par

\vspace{2.5pt plus 3pt minus 0.5pt}\noindent\hspace*{1.2em}\kbvnum{(೧)}\hspace{0.55em}ಕಬ್ಬು--ರಸ--ಸಿಹಿ\par

\vspace{2.5pt plus 3pt minus 0.5pt}\noindent\hspace*{1.2em}\kbvnum{(೨)}\hspace{0.55em}ರಾಮ--ರಾವಣ--ಸಂಹಾರ\par

\vspace{2.5pt plus 3pt minus 0.5pt}\noindent\hspace*{1.2em}\kbvnum{(೩)}\hspace{0.55em}ವ್ಯಾಸ--ಮಹಾಭಾರತ--ರಚನೆ\par

\vspace{2.5pt plus 3pt minus 0.5pt}\noindent\hspace*{1.2em}\kbvnum{(೪)}\hspace{0.55em}ಹಾಲು--ಮೊಸರು--ಬೆಣ್ಣೆ--ತುಪ್ಪ\par

\vspace{2.5pt plus 3pt minus 0.5pt}\noindent\hspace*{1.2em}\kbvnum{(೫)}\hspace{0.55em}ಐದು--ಮೂರು--ಎಂಟು\par

\kbvchapter{ಪಾಠ ೨}{ಅಕ್ಷರ ಮಾಲೆ}

\kbvsechead{ಕನ್ನಡ ಅಕ್ಷರ ಮಾಲೆ :}

\vspace{4pt plus 4pt minus 1pt}\begin{center}{\fontsize{11pt}{17pt}\selectfont \makebox[1.7em]{{\kbvsecfont ಅ}}\makebox[1.7em]{{\kbvsecfont ಆ}}\makebox[1.7em]{{\kbvsecfont ಇ}}\makebox[1.7em]{{\kbvsecfont ಈ}}\makebox[1.7em]{{\kbvsecfont ಉ}}\makebox[1.7em]{{\kbvsecfont ಊ}}\makebox[1.7em]{{\kbvsecfont ಋ}}\makebox[1.7em]{{\kbvsecfont ೠ}}\makebox[1.7em]{{\kbvsecfont ಎ}}\makebox[1.7em]{{\kbvsecfont ಏ}}\makebox[1.7em]{{\kbvsecfont ಐ}}\makebox[1.7em]{{\kbvsecfont ಒ}}\makebox[1.7em]{{\kbvsecfont ಓ}}\makebox[1.7em]{{\kbvsecfont ಔ}} \\
\makebox[1.7em]{ಂ}\makebox[1.7em]{(ಅಂ)}\makebox[1.7em]{;}\makebox[1.7em]{ಃ}\makebox[1.7em]{(ಅಃ)} \\
\makebox[1.7em]{ಕ}\makebox[1.7em]{ಖ}\makebox[1.7em]{ಗ}\makebox[1.7em]{ಘ}\makebox[1.7em]{ಙ} \\
\makebox[1.7em]{ಚ}\makebox[1.7em]{ಛ}\makebox[1.7em]{ಜ}\makebox[1.7em]{ಝ}\makebox[1.7em]{ಞ} \\
\makebox[1.7em]{ಟ}\makebox[1.7em]{ಠ}\makebox[1.7em]{ಡ}\makebox[1.7em]{ಢ}\makebox[1.7em]{ಣ} \\
\makebox[1.7em]{ತ}\makebox[1.7em]{ಥ}\makebox[1.7em]{ದ}\makebox[1.7em]{ಧ}\makebox[1.7em]{ನ} \\
\makebox[1.7em]{ಪ}\makebox[1.7em]{ಫ}\makebox[1.7em]{ಬ}\makebox[1.7em]{ಭ}\makebox[1.7em]{ಮ} \\
\makebox[1.7em]{ಯ}\makebox[1.7em]{ರ}\makebox[1.7em]{ಲ}\makebox[1.7em]{ವ}\makebox[1.7em]{ಶ}\makebox[1.7em]{ಷ}\makebox[1.7em]{ಸ}\makebox[1.7em]{ಹ}\makebox[1.7em]{ಳ}}\end{center}\vspace{4pt plus 4pt minus 1pt}

\begin{kbvrule}
\noindent\kbvbullet\hspace{0.5em}ಕನ್ನಡ ಅಕ್ಷರ ಮಾಲೆಯಲ್ಲಿರುವ ಬಿಡಿ ಅಕ್ಷರಗಳ ಸಂಖ್ಯೆ ಐವತ್ತು. ಅವುಗಳನ್ನು ಸ್ವರಗಳು, ವ್ಯಂಜನಗಳು ಮತ್ತು ಯೋಗವಾಹಗಳು ಎಂದು ಮೂರು ವಿಧವಾಗಿ ವಿಂಗಡಿಸಬಹುದು.
\end{kbvrule}

\kbvsechead{ಸ್ವರಗಳು :}

\begin{kbvrule}
\noindent\kbvbullet\hspace{0.5em}ಅ, ಆ, ಇ, ಈ, ಉ, ಊ, ಋ, ೠ, ಎ, ಏ, ಐ ಮತ್ತು ಒ, ಓ, ಔ ಸ್ವತಂತ್ರವಾಗಿ ಉಚ್ಚರಿಸಬಹುದಾದ ಅಕ್ಷರಗಳು ; ಇವುಗಳಿಗೆ ಸ್ವರಗಳೆಂದು ಹೆಸರು.
\end{kbvrule}
